\section{Aggregation, sparsity and fixed points}
\label{sec:theory}

We first analyse an exactly solvable orthogonal design, then characterise
fixed points for general designs. The orthogonal results distinguish
support enlargement caused by averaging from dependence on the number of
local epochs. They do not assume that every average of sparse vectors is
dense.

\subsection{Orthogonal designs}

Assume in this subsection that each site satisfies
$X_j^\top X_j=m_j I_p$, which requires $p\le\min_j m_j$, and write
$z_j=X_j^\top y_j/m_j$. Penalties are treated as fixed positive constants.

\begin{proposition}[Orthogonal designs]
\label{prop:orth}
Under these assumptions:
\begin{enumerate}
\item[(i)] For every integer $E\ge1$ and starting point $\bbeta\in\R^p$,
 $T_j^E(\bbeta)=\soft(z_j,\lambda_j)$.
\item[(ii)] The unrestricted sequence in Algorithm~\ref{alg:fedavg} satisfies
 $\bbeta_t=\bbar:=\sum_jw_j\soft(z_j,\lambda_j)$ for all $t\ge1$.
 With $\varepsilon_{\mathrm{tol}}>0$, the stopping rule terminates within
 two rounds. Its output does not depend on $E$.
\item[(iii)] The pooled design satisfies $X^\top X=mI_p$, and its unique
 minimiser is $\hat\bbeta_F=\soft(\bar z,\bar\lambda)$, where
 $\bar z=\sum_jw_jz_j$.
\item[(iv)] The supports satisfy
 $\supp(\hat\bbeta_F)\subseteq\bigcup_j\supp(\soft(z_j,\lambda_j))$.
 For fixed penalties and weights, this union equals $\supp(\bbar)$ except
 on a Lebesgue-null subset of $(z_1,\ldots,z_K)\in\R^{Kp}$ where non-zero
 contributions cancel exactly.
\end{enumerate}
\end{proposition}

\begin{proof}
For coordinate $k$, orthogonality gives
\[
 x_{jk}^\top(y_j-X_j\bbeta+x_{jk}\bbeta_k)
 =x_{jk}^\top y_j,
 \qquad \|x_{jk}\|_2^2=m_j.
\]
Equation~\eqref{eq:cd} therefore sets the coordinate to
$\soft(z_{jk},\lambda_j)$ independently of its starting value. One sweep
sets every coordinate to its local minimiser and subsequent sweeps leave
it unchanged. This proves (i). Consequently $\bbeta_1=\bbeta_2=\bbar$,
proving (ii); the stopping rule may already hold in round one.
For (iii), sum $X_j^\top X_j=m_jI_p$ over sites and apply
Lemma~\ref{lem:pooled} coordinatewise.

If $|z_{jk}|\le\lambda_j$ at every site, then
$|\bar z_k|\le\sum_jw_j|z_{jk}|\le\bar\lambda$, so coordinate $k$ is
absent from the pooled solution. This proves the inclusion in (iv).
To prove its second statement, partition $\R^{Kp}$ by the local active
sets and signs. On each of the finitely many resulting regions,
$\bbar_k$ is affine in the $z_{jk}$. If at least one site is active at
coordinate $k$, the coefficient of that site's $z_{jk}$ is $w_j>0$.
Thus $\bbar_k=0$ lies on a proper affine hyperplane. The union of these
hyperplanes over coordinates and regions, together with their boundaries,
has Lebesgue measure zero.
\end{proof}

The support inclusion need not be strict. For example, with $K=2$, equal
weights, $\lambda_1=\lambda_2=1$, $z_1=(2,0,0)$ and $z_2=(3,0,0)$,
the average and pooled solution both equal $(1.5,0,0)$. Their supports
remain equal and sparse on an open neighbourhood of these statistics.
Averaging can enlarge a support when sites select different coordinates;
it does not necessarily remove all exact zeros.

Part (ii) establishes exact independence from $E$ for orthogonal local
designs. Nonorthogonal designs can require multiple coordinate-descent
sweeps and hence can produce $E$-dependent iterates. This observation does
not provide a quantitative ordering across correlated designs. In
particular, independent predictors in the population do not yield
orthogonal empirical Gram matrices when $p>m_j$.

\begin{corollary}[False-positive probability under a fixed penalty]
\label{cor:fp}
In the setting of Proposition~\ref{prop:orth}, suppose $K\ge2$,
$\lambda_j=\lambda>0$, and, for a null coordinate $k$,
$z_{jk}\sim N(0,\sigma^2/m_j)$ independently across sites, with $\sigma>0$.
Then
\[
\Prob(k\in\supp(\bbar))
 =1-\prod_{j=1}^K\{1-2\bar\Phi(\lambda\sqrt{m_j}/\sigma)\}
 >2\bar\Phi(\lambda\sqrt m/\sigma)
 =\Prob(k\in\supp(\hat\bbeta_F)),
\]
where $\bar\Phi$ is the standard normal survival function. For fixed
$m$, $\lambda$ and $\sigma$, the left-hand side strictly increases with
$K$ along balanced partitions $m_j=m/K$ satisfying the design assumptions.
\end{corollary}

\begin{proof}
The Gaussian statistics have a joint density, so cancellation in
Proposition~\ref{prop:orth}(iv) has probability zero. Independence of the
events $\{|z_{jk}|>\lambda\}$ gives the first equality. Their union has
probability at least $2\bar\Phi(\lambda\sqrt{m_1}/\sigma)$, which strictly
exceeds $2\bar\Phi(\lambda\sqrt m/\sigma)$ because $m_1<m$.
Also $\bar z_k\sim N(0,\sigma^2/m)$, giving the pooled probability.
For balanced partitions let $q_K=2\bar\Phi(\lambda\sqrt{m/K}/\sigma)$.
Both $q_K\in(0,1)$ and $1-(1-q_K)^K$ strictly increase with $K$.
\end{proof}

The corollary concerns exact non-zero coefficients and a fixed common
penalty. It does not establish the same monotonicity for independently
tuned penalties, arbitrary repartitions, or coefficients declared active
only above a positive numerical threshold.

\subsection{General designs: fixed points}

Write $D_j^E(b)=T_j^E(b)-b$ for the local displacement. The following
statement characterises fixed points; it does not assert that the
averaged iteration converges to one.

\begin{proposition}[Averaged stationarity]
\label{prop:fixed}
A point $b^\dagger$ is a fixed point of Algorithm~\ref{alg:fedavg} if and
only if
\begin{equation}
 \sum_{j=1}^K w_jD_j^E(b^\dagger)=0.
 \label{eq:fixedpoint}
\end{equation}
All its local displacements vanish if and only if
$b^\dagger\in\bigcap_j\arg\min\Lloc_j$.
Thus, if this intersection is empty, any fixed point has non-zero local
displacements that cancel in the weighted average. In general,
\eqref{eq:fixedpoint} does not imply $0\in\partial F(b^\dagger)$.
\end{proposition}

\begin{proof}
The update equals $b+\sum_jw_jD_j^E(b)$, proving the first equivalence.
For a fixed site, each coordinate subproblem is strictly convex because
its column has positive squared norm. An exact coordinate update
therefore strictly decreases $\Lloc_j$ whenever it changes the coordinate.
If $T_j^E(b)=b$, the objective has the same value before and after the
sweeps, so no update can have changed a coordinate. The coordinatewise
optimality conditions then give
$0\in X_j^\top(X_jb-y_j)/m_j+\lambda_j\partial\|b\|_1$,
which is sufficient for a global minimum by convexity. Conversely, a
minimiser satisfies these conditions and every unique coordinate update
leaves it fixed. This proves the second equivalence.

For the final nonimplication take $K=2$, $p=m_1=m_2=1$,
$X_1=X_2=[1]$, $y_1=2$, $y_2=0$ and $\lambda_1=\lambda_2=1$.
For every $E\ge1$ the local maps return $1$ and $0$, respectively, so
$b^\dagger=1/2$ is an averaged fixed point. The pooled minimiser is
$\soft(1,1)=0$, while $F'(1/2)=1/2\ne0$.
\end{proof}

Distinct local minimisers do not by themselves exclude a pooled-optimal
fixed point; the sparse example following Proposition~\ref{prop:orth}
already shows this. The relevant diagnostic is the objective gap to a
numerically converged solution of the same pooled problem, rather than
the local displacements alone. Consensus ADMM provides that comparison
without requiring row-level exchange \citep{boyd2011}.

\begin{remark}
Proposition~\ref{prop:orth} requires $p\le\min_jm_j$ and does not apply
directly to our $p=600>m_j$ simulations. Proposition~\ref{prop:fixed}
applies to general designs but gives neither a convergence rate nor a
bound on the objective gap as a function of $E$. The experiments assess
support, optimisation accuracy and computational cost separately within
their specified designs.
\end{remark}
